\abstract{Portable and wireless EEG systems reduce setup burden and permit measurements closer to natural task conditions, but low channel count, motion, electrode technology, timing uncertainty, and opaque artifact handling can limit interpretation. This systematic scoping review will map wearable and low-density EEG studies in ADHD and pediatric neurodevelopmental research that use task-locked, intervention-locked, or state-monitoring paradigms. The review will extract device and electrode characteristics, sampling and event-marker methods, synchronization evidence, artifact detection/rejection, motion context, signal-quality reporting, task design, within-person versus group-level inference, reproducibility, and external validation. Existing reviews show both increasing use of wireless EEG in developmental populations and persistent limitations in data-quality reporting and artifact control. The present review therefore emphasizes measurement validity rather than diagnostic classification. Its intended contribution is a reporting and evidence framework for judging when wearable EEG supports reproducible event-locked inference in pediatric digital-intervention research.}
